\documentclass[11pt]{article}
\usepackage[utf8]{inputenc}
\usepackage[T1]{fontenc}
\usepackage{lmodern}
\usepackage{geometry}
\usepackage{amsmath,amssymb,amsthm,mathtools}
\usepackage{enumitem}
\usepackage{hyperref}
\usepackage{microtype}
\geometry{margin=1in}
\hypersetup{colorlinks=true,linkcolor=blue,urlcolor=blue,citecolor=blue}

\newcommand{\C}{\mathbb{C}}
\newcommand{\CP}{\mathbb{CP}}
\newcommand{\PP}{\mathbb{P}}
\newcommand{\OO}{\mathcal{O}}
\newcommand{\MM}{\mathcal{M}}
\newcommand{\Div}{\operatorname{Div}}
\newcommand{\Pic}{\operatorname{Pic}}
\newcommand{\ord}{\operatorname{ord}}
\newcommand{\Hom}{\operatorname{Hom}}
\newcommand{\im}{\operatorname{im}}
\newcommand{\rk}{\operatorname{rk}}
\newcommand{\Spec}{\operatorname{Spec}}
\newcommand{\degd}{\operatorname{deg}}
\newcommand{\ellD}{\ell}
\newcommand{\CPone}{\CP^1}
\newcommand{\into}{\hookrightarrow}
\newcommand{\surj}{\twoheadrightarrow}

\theoremstyle{definition}
\newtheorem{definition}{Definition}[section]
\newtheorem{example}[definition]{Example}
\newtheorem{exercise}{Exercise}[section]
\newtheorem{remark}[definition]{Remark}
\newtheorem{corollary}[definition]{Corollary}
\theoremstyle{plain}
\newtheorem{theorem}[definition]{Theorem}
\newtheorem{proposition}[definition]{Proposition}
\newtheorem{lemma}[definition]{Lemma}

\title{A Compute-First Lecture Note on Riemann--Roch for \texorpdfstring{$\CPone$}{CP1}\\[4pt]
with 30 Exercises}
\author{}
\date{}

\begin{document}
\maketitle
\tableofcontents

\section{Purpose of this note}
The goal of this note is to understand the Riemann--Roch theorem in the easiest nontrivial compact Riemann surface, namely
\[
X=\CPone.
\]
On a general compact Riemann surface, Riemann--Roch says
\[
\ell(D)-\ell(K-D)=\degd D+1-g,
\]
where
\begin{itemize}[leftmargin=2em]
    \item $D$ is a divisor,
    \item $K$ is a canonical divisor,
    \item $g$ is the genus,
    \item $\ell(D)=\dim_{\C} L(D)$, where
    \[
    L(D):=\{f\in \MM(X): (f)+D\ge 0\}\cup\{0\}.
    \]
\end{itemize}
For $X=\CPone$ we have $g=0$, so the theorem becomes
\[
\ell(D)-\ell(K-D)=\degd D+1.
\]
The miracle on $\CPone$ is that every divisor is linearly equivalent to a multiple of $[\infty]$, so the entire theorem reduces to explicit calculations with rational functions and polynomials.

The strategy of this note is:
\begin{enumerate}[label=\arabic*., leftmargin=2em]
    \item understand coordinates on $\CPone$,
    \item compute divisors of rational functions,
    \item show every divisor is equivalent to $n[\infty]$,
    \item compute $L(n[\infty])$ explicitly,
    \item compute the canonical divisor,
    \item verify Riemann--Roch directly.
\end{enumerate}

\section{The surface \texorpdfstring{$\CPone$}{CP1} and its coordinates}
Recall
\[
\CPone=(\C^2\setminus\{(0,0)\})/\C^\times,
\]
where
\[
\lambda\cdot (Z_0,Z_1)=(\lambda Z_0,\lambda Z_1),\qquad \lambda\in \C^\times.
\]
A point is written $[Z_0:Z_1]$.

There are two standard affine charts:
\[
U_0=\{Z_0\neq 0\},\qquad U_\infty=\{Z_1\neq 0\}.
\]
On $U_0$ we use the coordinate
\[
z=\frac{Z_1}{Z_0},
\]
so $U_0\cong \C$. On $U_\infty$ we use the coordinate
\[
w=\frac{Z_0}{Z_1},
\]
so $U_\infty\cong \C$. On the overlap $U_0\cap U_\infty$ we have
\[
w=\frac{1}{z}.
\]
Thus $\CPone$ is obtained by gluing two copies of $\C$ by inversion.

The point at infinity is
\[
\infty=[0:1],
\]
which lies in $U_\infty$ and corresponds to $w=0$.

\section{Meromorphic functions on \texorpdfstring{$\CPone$}{CP1}}
A meromorphic function on $\CPone$ is exactly a rational function in the affine coordinate $z$.

\begin{proposition}
We have
\[
\MM(\CPone)=\C(z).
\]
\end{proposition}

\begin{proof}
A rational function in $z$ is certainly meromorphic on the $z$-chart and also meromorphic near $\infty$ after the substitution $z=1/w$.
Conversely, let $f$ be meromorphic on $\CPone$. Restrict $f$ to $U_0\cong \C$; then $f$ is meromorphic on $\C$. Since $\infty$ is an isolated singularity of $f$, the function has at worst a pole at infinity. A meromorphic function on $\C$ with at worst a pole at infinity is rational. Hence $f\in \C(z)$.
\end{proof}

So every global meromorphic function is concrete: it is just a ratio $P(z)/Q(z)$ of polynomials.

\section{Divisors and orders on \texorpdfstring{$\CPone$}{CP1}}
\begin{definition}
A divisor on $X$ is a finite formal sum
\[
D=\sum_{p\in X} n_p[p],\qquad n_p\in \mathbb Z.
\]
Its degree is
\[
\degd D=\sum_{p\in X} n_p.
\]
\end{definition}

If $f$ is a nonzero meromorphic function and $p\in X$, then $\ord_p(f)$ denotes the vanishing order of $f$ at $p$ (negative if $f$ has a pole). The principal divisor of $f$ is
\[
(f)=\sum_{p\in X}\ord_p(f)[p].
\]

\begin{example}
In the coordinate $z$ on $U_0$, the function $z-a$ has a simple zero at $a\in \C$. To understand what happens at $\infty$, write $z=1/w$. Then
\[
z-a=\frac{1-aw}{w}.
\]
Near $w=0$, the factor $1-aw$ is a unit, while $1/w$ has a simple pole. Therefore
\[
(z-a)=[a]-[\infty].
\]
\end{example}

\begin{example}
More generally,
\[
\left(\frac{z-a}{z-b}\right)=[a]-[b].
\]
\end{example}

\begin{proposition}
Every principal divisor on $\CPone$ has degree $0$.
\end{proposition}

\begin{proof}
Let $f=P/Q$ with coprime polynomials. If $\deg P=m$ and $\deg Q=n$, then the finite zeros contribute $m$, the finite poles contribute $-n$, and near infinity we have
\[
f(1/w)=\frac{P(1/w)}{Q(1/w)}\sim c\,w^{n-m},
\]
so
\[
\ord_{\infty}(f)=n-m.
\]
Hence
\[
\degd (f)=m-n+(n-m)=0.
\]
\end{proof}

\section{Linear equivalence on \texorpdfstring{$\CPone$}{CP1}}
Two divisors $D,E$ are linearly equivalent, written $D\sim E$, if $D-E=(f)$ for some nonzero meromorphic function $f$.

The key fact on $\CPone$ is that degree is the only invariant of divisor class.

\begin{theorem}
Every divisor $D$ on $\CPone$ is linearly equivalent to $(\degd D)[\infty]$.
In particular,
\[
\Pic(\CPone)\cong \mathbb Z,
\qquad
D\longmapsto \degd D.
\]
\end{theorem}

\begin{proof}
Write
\[
D=\sum_{i=1}^r m_i[a_i]-\sum_{j=1}^s n_j[b_j]+c[\infty],
\]
where $a_i,b_j\in \C$ and $m_i,n_j>0$. Define
\[
h(z)=\frac{\prod_{i=1}^r (z-a_i)^{m_i}}{\prod_{j=1}^s (z-b_j)^{n_j}}.
\]
Then
\[
(h)=\sum_{i=1}^r m_i[a_i]-\sum_{j=1}^s n_j[b_j]-\left(\sum_{i=1}^r m_i-\sum_{j=1}^s n_j\right)[\infty].
\]
Therefore
\[
D-(h)=\left(c+\sum_{i=1}^r m_i-\sum_{j=1}^s n_j\right)[\infty]= (\degd D)[\infty].
\]
So $D\sim (\degd D)[\infty]$.
\end{proof}

This theorem is the concrete genus-zero form of the statement that every line bundle on $\CPone$ is some $\OO(n)$.

\section{The spaces \texorpdfstring{$L(D)$}{L(D)}}
For a divisor $D$ on $\CPone$, define
\[
L(D)=\{f\in \MM(\CPone):(f)+D\ge 0\}\cup\{0\}.
\]
Intuitively, $L(D)$ is the space of meromorphic functions whose poles are bounded by $D$.

Since every divisor is equivalent to $n[\infty]$, it suffices to compute $L(n[\infty])$.

\begin{proposition}
If $n<0$, then
\[
L(n[\infty])=0.
\]
If $n\ge 0$, then
\[
L(n[\infty])=\{\text{polynomials in }z\text{ of degree }\le n\}.
\]
In particular,
\[
\ellD(n[\infty])=
\begin{cases}
0,& n<0,\\
 n+1,& n\ge 0.
\end{cases}
\]
\end{proposition}

\begin{proof}
Assume first $n\ge 0$. If $f(z)=a_0+a_1z+\cdots +a_n z^n$ is a polynomial of degree at most $n$, then near $\infty$ we have
\[
f(1/w)=a_0+a_1w^{-1}+\cdots+a_n w^{-n},
\]
so $f$ has a pole of order at most $n$ at $\infty$ and no other poles. Hence $f\in L(n[\infty])$.

Conversely, let $f\in L(n[\infty])$. Then $f$ has no poles in $\C$ and a pole of order at most $n$ at $\infty$. Thus $f$ is holomorphic on $\C$, hence entire, and has polynomial growth of order at most $n$ at infinity. Therefore $f$ is a polynomial of degree at most $n$.

If $n<0$, then $f\in L(n[\infty])$ would have to vanish at infinity to order at least $-n$, while having no poles anywhere. So $f$ would be a holomorphic function on the compact surface $\CPone$ vanishing at a point. Since global holomorphic functions on $\CPone$ are constant, we get $f=0$.
\end{proof}

\begin{corollary}
For any divisor $D$ on $\CPone$,
\[
\ellD(D)=
\begin{cases}
0,& \degd D<0,\\
\degd D+1,& \degd D\ge 0.
\end{cases}
\]
\end{corollary}

\begin{proof}
Since $D\sim (\degd D)[\infty]$, the spaces $L(D)$ and $L((\degd D)[\infty])$ are naturally isomorphic by multiplication with a meromorphic function that realizes the equivalence. Hence dimensions agree, and the previous proposition applies.
\end{proof}

\section{The line bundles \texorpdfstring{$\OO(n)$}{O(n)} on \texorpdfstring{$\CPone$}{CP1}}
The standard line bundle $\OO(n)$ can be described by transition function $z^n$ on the overlap $U_0\cap U_\infty$.
A global section is a pair of holomorphic functions $(s_0,s_\infty)$ on the two charts satisfying
\[
s_\infty(w)=w^n s_0(1/w).
\]
If $n\ge 0$, then $s_0$ must be a polynomial of degree at most $n$.
Hence
\[
H^0(\CPone,\OO(n))\cong \operatorname{span}\{1,z,z^2,\dots,z^n\}
\]
as a vector space on the affine chart.
Equivalently, in homogeneous coordinates,
\[
H^0(\CPone,\OO(n))\cong \{
\text{homogeneous polynomials of degree }n\text{ in }Z_0,Z_1
\}.
\]
Thus
\[
h^0(\CPone,\OO(n))=
\begin{cases}
0,& n<0,\\
 n+1,& n\ge 0.
\end{cases}
\]

\begin{remark}
The divisor bundle of $n[\infty]$ is $\OO(n)$. So the explicit computation of $L(n[\infty])$ is exactly the explicit computation of global sections of $\OO(n)$.
\end{remark}

\section{The canonical divisor on \texorpdfstring{$\CPone$}{CP1}}
To apply Riemann--Roch we need the canonical divisor $K$.
On the chart $U_0$ the form $dz$ is holomorphic. Near infinity, write $z=1/w$. Then
\[
dz=d\left(\frac{1}{w}\right)=-\frac{1}{w^2}\,dw.
\]
Thus $dz$ has a pole of order $2$ at $w=0$, i.e. at $\infty$, and no zeros elsewhere. Therefore
\[
(dz)=-2[\infty].
\]
So a canonical divisor is
\[
K\sim -2[\infty],
\]
and hence
\[
K_{\CPone}\cong \OO(-2).
\]

\begin{corollary}
We have
\[
\ellD(K)=\ellD(-2[\infty])=0.
\]
In particular, $h^0(\CPone,K)=0$, as expected for genus $0$.
\end{corollary}

\section{Riemann--Roch on \texorpdfstring{$\CPone$}{CP1}}
Now we can verify Riemann--Roch directly.

\begin{theorem}[Riemann--Roch on $\CPone$]
For every divisor $D$ on $\CPone$,
\[
\ellD(D)-\ellD(K-D)=\degd D+1.
\]
Since $g(\CPone)=0$, this is exactly the general Riemann--Roch formula
\[
\ellD(D)-\ellD(K-D)=\degd D+1-g.
\]
\end{theorem}

\begin{proof}
Write $D\sim n[\infty]$, where $n=\degd D$. Then
\[
K-D\sim -2[\infty]-n[\infty]=-(n+2)[\infty].
\]
Hence
\[
\ellD(D)=\ellD(n[\infty]),
\qquad
\ellD(K-D)=\ellD(-(n+2)[\infty]).
\]
We now split into cases.

\smallskip
\noindent
\textbf{Case 1: $n\ge 0$.}
Then $\ellD(D)=n+1$. Also $-(n+2)<0$, so $\ellD(K-D)=0$. Thus
\[
\ellD(D)-\ellD(K-D)=(n+1)-0=n+1=\degd D+1.
\]

\smallskip
\noindent
\textbf{Case 2: $n=-1$.}
Then $\ellD(D)=0$ and $K-D\sim -[\infty]$, so $\ellD(K-D)=0$. Hence
\[
\ellD(D)-\ellD(K-D)=0=\degd D+1.
\]

\smallskip
\noindent
\textbf{Case 3: $n\le -2$.}
Then $\ellD(D)=0$, while
\[
K-D\sim -(n+2)[\infty],
\]
and now $-(n+2)\ge 0$. Therefore
\[
\ellD(K-D)=-(n+2)+1=-n-1.
\]
So
\[
\ellD(D)-\ellD(K-D)=0-(-n-1)=n+1=\degd D+1.
\]
This proves the theorem.
\end{proof}

\section{A useful reformulation with \texorpdfstring{$\OO(n)$}{O(n)}}
Since $K\cong \OO(-2)$, Riemann--Roch may be written as
\[
h^0(\CPone,\OO(n))-h^0(\CPone,\OO(-n-2))=n+1.
\]
Because
\[
h^0(\CPone,\OO(m))=
\begin{cases}
0,& m<0,\\
 m+1,& m\ge 0,
\end{cases}
\]
this identity is completely explicit.

By Serre duality one also gets
\[
h^1(\CPone,\OO(n))=h^0(\CPone,\OO(-n-2)).
\]
Hence
\[
h^1(\CPone,\OO(n))=
\begin{cases}
0,& n\ge -1,\\
-n-1,& n\le -2.
\end{cases}
\]

\section{What this teaches us}
The theorem on $\CPone$ is elementary enough to compute completely, but it already reveals the general philosophy of Riemann--Roch.

\begin{enumerate}[label=\arabic*., leftmargin=2em]
    \item \textbf{Divisors control allowed poles.} The space $L(D)$ measures meromorphic functions with pole bounds encoded by $D$.
    \item \textbf{Degree is the first approximation to dimension.} On $\CPone$, once $\degd D\ge 0$, the answer is exactly $\degd D+1$.
    \item \textbf{The canonical divisor measures the correction term.} On $\CPone$ the correction is especially simple because $K\sim -2[\infty]$.
    \item \textbf{Genus zero is rigid.} Every degree-zero divisor is principal, every line bundle is $\OO(n)$, and every global meromorphic function is rational.
\end{enumerate}

The general theorem for a compact Riemann surface of genus $g$ keeps exactly the same shape, but the term $\ellD(K-D)$ becomes a genuine global obstruction.

\section{Thirty exercises}
The exercises are arranged in six groups. The first groups are computational, and the later ones synthesize the whole story.

\subsection{Group A: coordinates and local orders}
\begin{exercise}
Show that the transition map between the two standard affine charts on $\CPone$ is
\[
w=\frac{1}{z}.
\]
Explain why $\infty$ corresponds to $w=0$.
\end{exercise}

\begin{exercise}
Compute
\[
\ord_{\infty}(z),\qquad \ord_{\infty}(1/z),\qquad \ord_{\infty}(z^m)
\]
for $m\in \mathbb Z$.
\end{exercise}

\begin{exercise}
Compute the divisor of $z-a$ for $a\in \C$.
\end{exercise}

\begin{exercise}
Compute the divisor of
\[
\frac{z-a}{z-b}
\]
for distinct $a,b\in \C$.
\end{exercise}

\begin{exercise}
Let
\[
P(z)=c\prod_{i=1}^r (z-a_i)^{m_i}
\]
with $c\neq 0$ and distinct $a_i$. Compute $(P)$ on $\CPone$.
\end{exercise}

\subsection{Group B: divisors and linear equivalence on \texorpdfstring{$\CPone$}{CP1}}
\begin{exercise}
Let
\[
f(z)=\frac{P(z)}{Q(z)}
\]
with coprime polynomials $P,Q$. Compute $(f)$ in terms of the zeros of $P$, the zeros of $Q$, and the order at $\infty$.
\end{exercise}

\begin{exercise}
Prove directly that every principal divisor on $\CPone$ has degree $0$.
\end{exercise}

\begin{exercise}
Let
\[
D=\sum_{i=1}^r m_i[a_i]-\sum_{j=1}^s n_j[b_j]
\]
with all $a_i,b_j\in \C$ and $\degd D=0$. Construct explicitly a rational function $h$ with $(h)=D$.
\end{exercise}

\begin{exercise}
Let $D$ be any divisor on $\CPone$. Prove that
\[
D\sim (\degd D)[\infty].
\]
\end{exercise}

\begin{exercise}
Deduce that every degree-zero divisor on $\CPone$ is principal.
\end{exercise}

\subsection{Group C: the spaces \texorpdfstring{$L(D)$}{L(D)}}
\begin{exercise}
Prove that if $n<0$, then
\[
L(n[\infty])=0.
\]
\end{exercise}

\begin{exercise}
Prove that if $n\ge 0$, then
\[
L(n[\infty])=\operatorname{span}_{\C}\{1,z,z^2,\dots,z^n\}.
\]
\end{exercise}

\begin{exercise}
Compute $\ellD(n[\infty])$ for all integers $n$.
\end{exercise}

\begin{exercise}
Let $D=[a_1]+\cdots+[a_n]$ with $a_i\in \C$ not necessarily distinct. Show that
\[
\ellD(D)=n+1.
\]
Hint: first show $D\sim n[\infty]$.
\end{exercise}

\begin{exercise}
Let $D$ be a divisor on $\CPone$ of degree $d$. Show that
\[
\ellD(D)=0\quad\text{if }d<0,
\qquad
\ellD(D)=d+1\quad\text{if }d\ge 0.
\]
\end{exercise}

\subsection{Group D: the bundles \texorpdfstring{$\OO(n)$}{O(n)}}
\begin{exercise}
Using the transition function $z^n$ on $U_0\cap U_\infty$, show that a global section of $\OO(n)$ is given by a pair $(s_0,s_\infty)$ satisfying
\[
s_\infty(w)=w^n s_0(1/w).
\]
\end{exercise}

\begin{exercise}
Show that for $n\ge 0$,
\[
H^0(\CPone,\OO(n))\cong \operatorname{span}_{\C}\{1,z,\dots,z^n\}
\]
on the affine chart $U_0$.
\end{exercise}

\begin{exercise}
Show that for $n\ge 0$,
\[
H^0(\CPone,\OO(n))
\]
may also be identified with homogeneous polynomials of degree $n$ in $Z_0,Z_1$.
\end{exercise}

\begin{exercise}
Compute
\[
h^0(\CPone,\OO(n))
\]
for all integers $n$.
\end{exercise}

\begin{exercise}
Show that if $n<0$, then $\OO(n)$ has no nonzero global holomorphic sections.
\end{exercise}

\subsection{Group E: canonical divisor and Riemann--Roch}
\begin{exercise}
Compute the divisor of $dz$ on $\CPone$ and show that
\[
K\sim -2[\infty].
\]
\end{exercise}

\begin{exercise}
Compute the divisor of $z^m dz$ for an integer $m$.
\end{exercise}

\begin{exercise}
Show that
\[
\ellD(K)=0.
\]
Interpret this as $h^0(\CPone,K)=0$.
\end{exercise}

\begin{exercise}
Verify Riemann--Roch directly for the divisor $D=n[\infty]$ for every integer $n$.
\end{exercise}

\begin{exercise}
Deduce from the previous exercise that Riemann--Roch holds for every divisor $D$ on $\CPone$.
\end{exercise}

\subsection{Group F: synthesis and applications}
\begin{exercise}
Use Riemann--Roch and the formula $K\sim -2[\infty]$ to compute
\[
h^1(\CPone,\OO(n))
\]
for every integer $n$.
\end{exercise}

\begin{exercise}
Show that the degree map induces an isomorphism
\[
\Pic(\CPone)\cong \mathbb Z,
\qquad
\OO(n)\longmapsto n.
\]
\end{exercise}

\begin{exercise}
Let $D$ be an effective divisor of degree $n$ on $\CPone$. Show that there exists a polynomial $P(z)$ of degree exactly $n$ such that the zero divisor of $P$ on $\C$ is exactly $D-n[\infty]$.
\end{exercise}

\begin{exercise}
Show that the complete linear system $|n[\infty]|$ defines a holomorphic map
\[
\CPone\longrightarrow \PP^n,
\qquad
[z_0:z_1]\longmapsto [z_0^n:z_0^{n-1}z_1:\dots:z_1^n].
\]
Describe its image.
\end{exercise}

\begin{exercise}
Explain in your own words why $\CPone$ is the perfect first test case for Riemann--Roch. Your explanation should mention at least the following points:
\begin{enumerate}[label=(\alph*), leftmargin=2em]
    \item every meromorphic function is rational,
    \item every divisor class is determined by degree,
    \item $K\sim -2[\infty]$,
    \item the dimension formula becomes completely explicit.
\end{enumerate}
\end{exercise}

\section{Hints and brief solutions}
\subsection*{1--5}
\begin{enumerate}[label=\textbf{\arabic*.}, leftmargin=2.5em]
    \item By definition $z=Z_1/Z_0$ on $U_0$ and $w=Z_0/Z_1$ on $U_\infty$, hence $w=1/z$. The point $\infty=[0:1]$ satisfies $w=0$.
    \item Since $z=1/w$ near $\infty$, we get $\ord_{\infty}(z)=-1$, $\ord_{\infty}(1/z)=1$, and $\ord_{\infty}(z^m)=-m$.
    \item One zero at $z=a$ and one pole at $\infty$, so $(z-a)=[a]-[\infty]$.
    \item Subtract the previous divisors: $((z-a)/(z-b))=[a]-[b]$.
    \item The finite zeros contribute $\sum m_i$, and at infinity there is a pole of order $\sum m_i$. Thus
    \[
    (P)=\sum_{i=1}^r m_i[a_i]-\left(\sum_{i=1}^r m_i\right)[\infty].
    \]
\end{enumerate}

\subsection*{6--10}
\begin{enumerate}[label=\textbf{\arabic*.}, leftmargin=2.5em, start=6]
    \item Write the finite zeros and poles explicitly. At infinity,
    \[
    \ord_{\infty}(P/Q)=\deg Q-\deg P.
    \]
    \item Sum the finite orders and the order at infinity; they cancel.
    \item Take
    \[
    h(z)=\frac{\prod (z-a_i)^{m_i}}{\prod (z-b_j)^{n_j}}.
    \]
    Degree zero guarantees no extra order at infinity.
    \item Construct $h$ from all finite points as in the proof in the note. Then $D-(h)$ is supported only at $\infty$, and its coefficient is $\degd D$.
    \item If $\degd D=0$, then the previous exercise gives $D\sim 0$, hence $D=(h)$.
\end{enumerate}

\subsection*{11--15}
\begin{enumerate}[label=\textbf{\arabic*.}, leftmargin=2.5em, start=11]
    \item A function in $L(n[\infty])$ would have no poles anywhere and would have to vanish at infinity. Global holomorphic functions on $\CPone$ are constant, so only $0$ works.
    \item A function in $L(n[\infty])$ is entire on $\C$ and has pole order at most $n$ at infinity, hence is a polynomial of degree at most $n$.
    \item The answer is $0$ for $n<0$ and $n+1$ for $n\ge 0$.
    \item Since each $[a_i]\sim [\infty]$, we get $D\sim n[\infty]$, so $\ellD(D)=\ellD(n[\infty])=n+1$.
    \item Use $D\sim d[\infty]$ where $d=\degd D$.
\end{enumerate}

\subsection*{16--20}
\begin{enumerate}[label=\textbf{\arabic*.}, leftmargin=2.5em, start=16]
    \item This is exactly the cocycle condition for sections of a line bundle with transition function $z^n$.
    \item Holomorphicity of $s_\infty(w)=w^n s_0(1/w)$ at $w=0$ forces $s_0$ to be a polynomial of degree at most $n$.
    \item A homogeneous degree-$n$ polynomial
    \[
    F(Z_0,Z_1)=\sum_{k=0}^n a_k Z_0^{n-k}Z_1^k
    \]
    becomes the polynomial $\sum a_k z^k$ on $U_0$.
    \item Thus $h^0(\CPone,\OO(n))=0$ for $n<0$ and $n+1$ for $n\ge 0$.
    \item If a nonzero global section existed for $n<0$, it would correspond to a divisor of negative degree but effective divisor of zeros, impossible.
\end{enumerate}

\subsection*{21--25}
\begin{enumerate}[label=\textbf{\arabic*.}, leftmargin=2.5em, start=21]
    \item Use $z=1/w$ and compute
    \[
    dz=-w^{-2}\,dw.
    \]
    Hence $(dz)=-2[\infty]$.
    \item Combine $(z^m)=m[0]-m[\infty]$ with $(dz)=-2[\infty]$ to get
    \[
    (z^m dz)=m[0]-(m+2)[\infty].
    \]
    \item Since $K\sim -2[\infty]$ has negative degree, $\ellD(K)=0$.
    \item Do the three cases $n\ge 0$, $n=-1$, and $n\le -2$.
    \item Since every divisor is equivalent to $n[\infty]$, the previous exercise is enough.
\end{enumerate}

\subsection*{26--30}
\begin{enumerate}[label=\textbf{\arabic*.}, leftmargin=2.5em, start=26]
    \item Use
    \[
    h^0(\OO(n))-h^1(\OO(n))=n+1
    \]
    together with $h^1(\OO(n))=h^0(\OO(-n-2))$ or solve case by case.
    \item Every line bundle is $\OO(D)$ for some divisor $D$, and $D\sim (\degd D)[\infty]$. Thus $\OO(D)\cong \OO(\degd D)$.
    \item Take the monic polynomial
    \[
    P(z)=\prod (z-a_i)^{m_i}
    \]
    where $D=\sum m_i[a_i]$ on the affine chart; its divisor is $D-n[\infty]$.
    \item The image is the rational normal curve of degree $n$ in $\PP^n$.
    \item A good answer should stress that on $\CPone$ every ingredient of Riemann--Roch can be computed by hand.
\end{enumerate}

\section{Final summary}
For $\CPone$ the Riemann--Roch theorem is not mysterious at all. Everything boils down to four explicit facts:
\begin{enumerate}[label=\arabic*., leftmargin=2em]
    \item $\MM(\CPone)=\C(z)$,
    \item every divisor satisfies $D\sim (\degd D)[\infty]$,
    \item $L(n[\infty])$ is the space of polynomials of degree at most $n$,
    \item $K\sim -2[\infty]$.
\end{enumerate}
Therefore
\[
\ellD(D)-\ellD(K-D)=\degd D+1
\]
becomes a direct polynomial count.

This is exactly why the genus-zero case should be mastered first: it turns the general theorem into a sequence of concrete computations with rational functions, divisors, and the point at infinity.

\end{document}
